\section{Distinct Subsequences}
\label{sec:dp-distinct-subsequences}

\problemstatement{Given strings $s$ (length $n$) and $t$ (length $m$),
count the number of distinct ways $t$ can be formed as a subsequence of
$s$ (counting different choices of \emph{positions} in $s$ as distinct,
even if they spell the same characters).}

\begin{patternbox}
\emph{``Count the number of ways one string appears as a subsequence of
another''} reuses Section~\ref{sec:dp-lcs}'s match/mismatch branching,
but on a \textbf{match}, there are now \textbf{two} valid choices to
\emph{sum} rather than one branch to take unconditionally: use the
current character of $s$ to match $t$, \emph{or} skip it and hope a
later character of $s$ matches instead -- both contribute to the count,
since $s$ may contain the needed character multiple times.
\end{patternbox}

\subsection*{Deriving the recurrence}

Let $f(i,j)$ be the number of ways $t[0..j-1]$ (length $j$) can be
formed as a subsequence of $s[0..i-1]$ (length $i$).

\[
  f(i,j) =
  \begin{cases}
    f(i{-}1,j{-}1) + f(i{-}1,j) & \text{if } s[i{-}1]=t[j{-}1], \\
    f(i{-}1,j) & \text{otherwise.}
  \end{cases}
\]

If the characters match, we may either \textbf{use} this occurrence in
$s$ to match $t$'s current character (advancing both, $f(i{-}1,j{-}1)$),
\emph{or} \textbf{skip} it, leaving $t$'s current character to be
matched by some earlier-yet-unconsidered part of $s$ ($f(i{-}1,j)$) --
both possibilities lead to valid, distinct subsequence selections, so
their counts are summed. If the characters don't match, this position
in $s$ cannot contribute to matching $t[j-1]$ at all, so only the skip
option remains.

Base cases: $f(i,0)=1$ for every $i$ (the empty string is a subsequence
of anything, in exactly one way -- selecting nothing). $f(0,j)=0$ for
every $j>0$ (a non-empty $t$ cannot be a subsequence of an empty $s$).

\subsection*{Approach 1: Recursion}

\begin{lstlisting}
#include <bits/stdc++.h>
using namespace std;

int distinctSubsequencesRecursive(int i, int j, string& s, string& t) {
    if (j == 0) return 1;
    if (i == 0) return 0;

    if (s[i - 1] == t[j - 1]) {
        return distinctSubsequencesRecursive(i - 1, j - 1, s, t) +
               distinctSubsequencesRecursive(i - 1, j, s, t);
    }
    return distinctSubsequencesRecursive(i - 1, j, s, t);
}

int main() {
    string s = "babgbag", t = "bag";
    cout << distinctSubsequencesRecursive(s.size(), t.size(), s, t) << endl; // 5
    return 0;
}
\end{lstlisting}

\begin{complexitybox}[Approach 1 Complexity]
\textbf{Time.} $O(2^n)$ in the worst case.
\textbf{Space.} $O(n)$ recursion stack.
\end{complexitybox}

\subsection*{Approach 2: Memoization}

\begin{lstlisting}
#include <bits/stdc++.h>
using namespace std;

int distinctSubsequencesMemo(int i, int j, string& s, string& t,
                              vector<vector<long long>>& memo) {
    if (j == 0) return 1;
    if (i == 0) return 0;
    if (memo[i][j] != -1) return memo[i][j];

    long long result;
    if (s[i - 1] == t[j - 1]) {
        result = distinctSubsequencesMemo(i - 1, j - 1, s, t, memo) +
                 distinctSubsequencesMemo(i - 1, j, s, t, memo);
    } else {
        result = distinctSubsequencesMemo(i - 1, j, s, t, memo);
    }
    return memo[i][j] = result;
}

int main() {
    string s = "babgbag", t = "bag";
    int n = s.size(), m = t.size();
    vector<vector<long long>> memo(n + 1, vector<long long>(m + 1, -1));
    cout << distinctSubsequencesMemo(n, m, s, t, memo) << endl; // 5
    return 0;
}
\end{lstlisting}

\begin{complexitybox}[Approach 2 Complexity]
\textbf{Time.} $O(n \cdot m)$.
\textbf{Space.} $O(n \cdot m)$ memo $+$ $O(n)$ recursion stack.
\end{complexitybox}

\subsection*{Approach 3: Tabulation}

\begin{lstlisting}
#include <bits/stdc++.h>
using namespace std;

long long distinctSubsequencesTabulation(string& s, string& t) {
    int n = s.size(), m = t.size();
    vector<vector<long long>> dp(n + 1, vector<long long>(m + 1, 0));

    for (int i = 0; i <= n; i++) dp[i][0] = 1;

    for (int i = 1; i <= n; i++) {
        for (int j = 1; j <= m; j++) {
            if (s[i - 1] == t[j - 1]) {
                dp[i][j] = dp[i - 1][j - 1] + dp[i - 1][j];
            } else {
                dp[i][j] = dp[i - 1][j];
            }
        }
    }
    return dp[n][m];
}

int main() {
    string s = "babgbag", t = "bag";
    cout << distinctSubsequencesTabulation(s, t) << endl; // 5
    return 0;
}
\end{lstlisting}

\subsection*{Dry run of the tabulation table}

\dryrun{Take $s=\texttt{"babgbag"}$, $t=\texttt{"bag"}$.}

The five ways \texttt{"bag"} appears as a subsequence of
\texttt{"babgbag"} correspond to choosing one of the two \texttt{'b'}s
(positions $0,2,4$... specifically $3$ valid \texttt{'b'} positions
combined with valid later \texttt{'a'},\texttt{'g'} pairs) in
combination with the available \texttt{'a'} and \texttt{'g'}
occurrences after it; building the table row by row accumulates these
combinations automatically via the sum rule, reaching
$\textit{dp}[7][3]=5$.

\begin{complexitybox}[Approach 3 Complexity]
\textbf{Time.} $O(n \cdot m)$.
\textbf{Space.} $O(n \cdot m)$ table, no recursion stack.
\end{complexitybox}

\subsection*{Approach 4: Space Optimization}

\begin{lstlisting}
#include <bits/stdc++.h>
using namespace std;

long long distinctSubsequencesSpaceOptimized(string& s, string& t) {
    int n = s.size(), m = t.size();
    vector<long long> prev(m + 1, 0), current(m + 1, 0);
    prev[0] = current[0] = 1;

    for (int i = 1; i <= n; i++) {
        for (int j = 1; j <= m; j++) {
            if (s[i - 1] == t[j - 1]) {
                current[j] = prev[j - 1] + prev[j];
            } else {
                current[j] = prev[j];
            }
        }
        prev = current;
    }
    return prev[m];
}

int main() {
    string s = "babgbag", t = "bag";
    cout << distinctSubsequencesSpaceOptimized(s, t) << endl; // 5
    return 0;
}
\end{lstlisting}

\begin{complexitybox}[Approach 4 Complexity]
\textbf{Time.} $O(n \cdot m)$.
\textbf{Space.} $O(m)$ for two rolling rows.
\end{complexitybox}

\begin{takeawaybox}
A character match does not always mean ``take it unconditionally and
advance both pointers'' -- when the same character could plausibly be
matched by \emph{several} different positions in the longer string (as
here, where $s$ may contain repeats), a match should branch into
\textbf{both} ``use it'' and ``save it for later, skip for now,''
summing their counts. This is a subtler version of the match case than
Section~\ref{sec:dp-lcs}'s, worth contrasting carefully.
\end{takeawaybox}
